% ============================================================
%  Chapter 10 — Stay Safe and Smart
%  Inside a Smartphone: The Tiny World in Your Hand
% ============================================================

\chapter{Stay Safe and Smart}

\chapteropening{%
  You have explored every part of a smartphone.
  Now it is time to learn how to use this powerful tool wisely,
  safely, and kindly.
}
\chapterbadge{Final Mission: Become a Smart Smartphone Explorer!}

% --- Opening ---
\section*{The Smartphone Explorer's Code}

All of the characters gathered around Sam one last time:
Chip, Memo, Tappy, Lens, Zap, Sensa, Appy, Orbi — and of course, Dot.

\character{Dot}{You have come a long way since that first question:
  ``How does all of this fit inside one tiny phone?''
  Now you know what is inside the tiny computer in your pocket.
  But knowing how something works also means knowing how to use it responsibly.
  Are you ready for your final challenge?}

\sectionrule

% --- Passwords ---
\section*{Tool 1: A Strong Password}

A \term{password} protects your accounts from people who should not access them.
A strong password:
\begin{itemize}
  \item Is long (at least 12 characters).
  \item Mixes letters, numbers, and symbols.
  \item Is different for every account.
  \item Is never shared with anyone except a trusted adult.
\end{itemize}

\begin{picturethis}
  A weak password is like a front door left wide open.
  A strong password is a sturdy, locked door.
  And sharing your password with someone you do not trust
  is like giving a stranger the key to your home.
\end{picturethis}

\sectionrule

% --- Biometrics ---
\section*{Tool 2: Fingerprints and Face Unlock}

\term{Biometrics} uses a unique part of your body to identify you.
Smartphones can unlock using:
\begin{itemize}
  \item \textbf{Fingerprint}: A sensor scans the unique ridges on your fingertip.
  \item \textbf{Face unlock}: Cameras and sensors map the unique shape of your face.
\end{itemize}

Biometrics are convenient because you always have them with you —
but a password is still a good backup.
If biometrics ever fail, your password helps you get back into your device safely.

\sectionrule

% --- App permissions ---
\section*{Tool 3: Asking Permission Before Downloading Apps}

\begin{picturethis}
  Imagine you invite a new friend to your house.
  They walk in and immediately start opening every cupboard, every drawer, and every door.
  You would say: ``Wait — you need to ask first!''

  Apps are the same. When a new app arrives on your phone,
  it should ask permission before looking at your photos, listening through
  the microphone, or checking where you are.
\end{picturethis}

When you install an app, it may ask for \term{permissions} —
access to parts of your phone like the camera, microphone, or location.

\begin{itemize}
  \item Always read what an app is asking permission for.
  \item Ask yourself: does this app really need this access?
        (A torch app does not need your location.)
  \item Ask a trusted adult before installing unfamiliar apps.
  \item Only download from trusted app stores.
\end{itemize}

\begin{quickmap}
  \textbf{Smart permission choices:}
  \begin{center}
  \begin{tabular}{p{0.6\linewidth}l}
    \textbf{Situation} & \textbf{Smart?} \\
    \hline
    Maps app asks for your location.          & Yes! \\
    A game asks to read all your contacts.    & Ask an adult. \\
    A torch app asks to use your microphone.  & Suspicious! \\
    You update your apps regularly.           & Great idea! \\
  \end{tabular}
  \end{center}
\end{quickmap}

\sectionrule

% --- Kind online ---
\section*{Tool 4: Being Kind Online}

A smartphone connects you to real people all over the world.
\begin{itemize}
  \item Treat people online the same way you would face to face.
  \item Think before you send: \emph{``Would I say this in person?''}
  \item If someone is unkind to you online, it is not your fault. Tell a trusted adult.
\end{itemize}

\character{Appy}{Messages sent through me reach real people with real feelings.
  A kind message can make someone's day.
  Think before you type.}

\sectionrule

% --- Screen time ---
\section*{Tool 5: Balance Your Screen Time}

\begin{picturethis}
  Sam looked at the clock and realised three hours had passed — it felt like twenty minutes!
  That is how smartphones work sometimes: time seems to disappear.
  Chip, Appy, and Tappy all work so well together that using the phone feels easy and fun.
  But the best explorers know when to put the map away and look up at the real sky.
\end{picturethis}

Smartphones are powerful tools — but so is the world away from screens.
\begin{itemize}
  \item Take regular breaks from your phone.
  \item Notice how you feel after a long session.
  \item Balance phone time with outdoor play, reading, and time with family and friends.
  \item Ask yourself: ``Am I choosing to use my phone right now, or just using it out of habit?''
\end{itemize}

\sectionrule

% --- Code of honour ---
\section*{Smartphone Explorer's Code of Honour}

\begin{picturethis}
  A great explorer does not just learn facts — they make promises.
  Read this code and sign it as a true smartphone explorer.
\end{picturethis}

\begin{enumerate}[label=\textbf{\arabic*.}]
  \item I will protect my passwords and keep personal information private.
  \item I will check app permissions carefully before installing anything new.
  \item I will keep my apps and operating system updated.
  \item I will be kind in every message, comment, and online game.
  \item I will ask a trusted adult for help if something feels wrong.
  \item I will balance screen time with rest, play, and real-world adventures.
\end{enumerate}

\needspace{6\baselineskip}
\vspace{6pt}
\begin{center}
  \begin{tcolorbox}[colback=SunYellow!15,colframe=SunYellow!70!black,arc=8pt,boxrule=1pt,width=0.9\linewidth]
    \textbf{Explorer's Signature:} \rule{0.48\linewidth}{0.4pt}\hfill
    \textbf{Date:} \rule{0.2\linewidth}{0.4pt}
  \end{tcolorbox}
\end{center}

\sectionrule

% --- Farewell ---
\section*{You Are a Smartphone Explorer!}

All the characters gathered around Sam and cheered.

\character{Dot}{You started this journey wondering what was inside a tiny phone.
  Now you know about the processor, memory, touchscreen, cameras, battery,
  sensors, apps, and GPS.
  You understand how they work together — and how to use them wisely.
  Take that knowledge and go explore the world!}

\begin{bigidea}
  A \textbf{smart smartphone user} understands their device,
  protects their information, treats others with kindness,
  and balances screen time with the rest of life.
\end{bigidea}
